% TOP_PROBLEM: {"id": 1484, "title": "Unknotting Problem", "category_id": 7, "category": {"id": 7, "name": "topology", "display_name": "Topology", "description": "Properties preserved under continuous deformations.", "slug": "topology", "order_index": 7, "created_at": "2026-07-31T15:26:25.671Z"}, "status": "open"}
% FIRST_POSTED: null
% ENTRY_KIND: statement_only
% Imported from: batch08.tex; UnsolvedMath ID 1484
\documentclass[11pt]{article}
\usepackage[margin=25mm]{geometry}
\usepackage{fontspec}
\setmainfont{Latin Modern Roman}
\usepackage{amsmath,amssymb,mathtools}
\usepackage{unicode-math}
\setmathfont{Latin Modern Math}
\usepackage{xurl,hyperref}
\hypersetup{colorlinks=true,urlcolor=blue}
\setcounter{secnumdepth}{0}
\setlength{\parindent}{0pt}
\setlength{\parskip}{5pt}
\setlength{\emergencystretch}{3em}
\newcommand{\sref}[2]{\hyperlink{source:#1:#2}{[S#2]}}
\newcommand{\cataloguescope}[1]{\paragraph{Recorded target.}#1}
\begin{document}
% UnsolvedMath ID 1484; source catalogue code TOP-006
\setcounter{section}{1483}
\section{Unknotting Problem}\label{1484}
{\small\sffamily UnsolvedMath ID 1484 · Topology}
\subsection{Definitions and mathematical statement}

\paragraph{Shared notation.}$\mathbb N=\{1,2,\ldots\}$, and $\mathbb Z,\mathbb Q,\mathbb R,\mathbb C$ denote the integers, rationals, reals and complex numbers. Any other conventions are specified locally.


A knot is a tame embedding of a circle in \(S^3\); equivalently it may
be represented by a piecewise-linear simple closed curve.
A knot diagram is a generic planar projection, with an over/under choice
at every crossing. Encode its finite planar graph, cyclic ordering of
incident arcs and crossing choices by a binary string of length \(L\).
A diagram represents the unknot if its knot is ambient isotopic to a
round circle: an ambient isotopy is a continuous deformation through
homeomorphisms of \(S^3\). Polynomial time means worst-case time
at most \(C(1+L)^c\), for constants \(C,c\) independent of the input,
on a classical deterministic Turing machine.
\paragraph{Statement recorded in the supplied source.}
Is there a polynomial-time algorithm which, on every valid
one-component knot diagram, decides whether that diagram represents
the unknot? \sref{1484}{2}
\subsection{Short English statement}
Can a classical algorithm decide whether any given knot diagram is an unknot in time bounded by a fixed polynomial in the input length?
\subsection{Sources}
\begin{description}
\item[\hypertarget{source:1484:1}{S1}] ULAM.AI and the UnsolvedMath Contributors, \emph{UnsolvedMath: A Curated Collection of Open Mathematics Problems}, record 1484 (TOP-006). CC BY 4.0. \url{https://huggingface.co/datasets/ulamai/UnsolvedMath}.
\item[\hypertarget{source:1484:2}{S2}] Joel Hass, Jeffrey C. Lagarias and Nicholas Pippenger, \emph{The Computational Complexity of Knot and Link Problems}, Journal of the ACM 46 (1999), 185--211; arXiv:math/9807016, Sections 1 and 3, knot diagrams and their input length. \url{https://arxiv.org/abs/math/9807016}.
\end{description}
\label{end:1484}
\end{document}
